\section*{Tóm tắt đề tài}
\phantomsection
\addcontentsline{toc}{section}{Tóm tắt đề tài}
Trong bối cảnh phương tiện giao thông điện, đặc biệt là xe máy điện, đang phát triển mạnh mẽ nhằm hướng tới mục tiêu giảm thiểu phát thải và bảo vệ môi trường, nhu cầu về một hệ thống hạ tầng sạc thông minh và tiện lợi ngày càng trở nên cấp thiết. Tuy nhiên, người sử dụng xe điện hiện nay thường xuyên gặp phải tình trạng khó khăn trong việc tìm kiếm trạm sạc, không thể biết trước tình trạng khả dụng của cổng sạc, dẫn đến lãng phí thời gian chờ đợi.

Để giải quyết vấn đề trên, đồ án nghiên cứu và phát triển hệ thống \textbf{EVGo}, một nền tảng quản lý và hỗ trợ đặt chỗ trạm sạc xe máy điện. Hệ thống được thiết kế theo kiến trúc Modular Monolith dựa trên Spring Boot (Spring Modulith) ở phía Backend, kết hợp với ứng dụng di động đa nền tảng React Native (Expo) dành cho người dùng cuối và giao diện quản trị Web dành cho chủ trạm sạc. 

Các chức năng cốt lõi của hệ thống bao gồm:
\begin{itemize}
    \item Tìm kiếm trạm sạc trên bản đồ số với thông tin hiển thị theo thời gian thực.
    \item Đặt chỗ cổng sạc trước thông qua cơ chế khóa phân tán giúp chống trùng lịch.
    \item Giám sát quá trình sạc theo thời gian thực nhờ tích hợp luồng dữ liệu Server-Sent Events và giao thức điều khiển trạm sạc tiêu chuẩn OCPP.
    \item Thanh toán trực tuyến tự động thông qua tích hợp cổng thanh toán ZaloPay.
    \item Cung cấp giao diện quản trị cho phép chủ trạm sạc quản lý trạm sạc của mình.
\end{itemize}

Kết quả của đồ án là một sản phẩm phần mềm hoàn chỉnh, hoạt động ổn định, đáp ứng được các luồng nghiệp vụ phức tạp từ khi người dùng bắt đầu tìm kiếm trạm sạc đến khi hoàn tất quá trình thanh toán. Hệ thống không chỉ mang lại trải nghiệm thuận tiện, minh bạch cho người sử dụng xe điện mà còn cung cấp công cụ vận hành hiệu quả cho các cá nhân, tổ chức muốn vận hành và tham gia kinh doanh hạ tầng trạm sạc.
